\documentclass[pmlr]{jmlr}

\usepackage{booktabs}
\usepackage{siunitx}
\usepackage{multirow}
\usepackage{graphicx}
\usepackage[switch]{lineno}
\usepackage{tikz}
\usepackage{amsmath}
\usepackage{amssymb}
\usepackage{algorithm}
\usepackage{algorithmic}
\usepackage{colortbl}
\usepackage{array}
\usepackage{hyperref}
\usetikzlibrary{arrows.meta, positioning, shapes.geometric, fit,
                backgrounds, calc, shadows, matrix, decorations.pathreplacing}
\DeclareMathOperator*{\argmax}{arg\,max}

% ── Color palette ─────────────────────────────────────────────────────────────
\definecolor{colonog}{RGB}{25,60,110}
\definecolor{ctspine}{RGB}{120,42,0}
\definecolor{ctpelv}{RGB}{80,22,100}
\definecolor{patdb}{RGB}{18,78,40}
\definecolor{placed}{RGB}{160,20,30}
\definecolor{tabhead}{RGB}{40,48,62}
\definecolor{rowalt}{RGB}{248,250,255}
\definecolor{winner}{RGB}{14,120,55}
\definecolor{fuseblue}{RGB}{25,100,180}
\definecolor{sepamber}{RGB}{185,100,0}
\definecolor{arrowgray}{RGB}{80,80,90}
\definecolor{panelbg}{RGB}{247,249,252}
\definecolor{metalgray}{RGB}{95,105,120}
\definecolor{lstvcol}{RGB}{180,30,30}
\definecolor{agreem}{RGB}{14,120,55}
\definecolor{disagr}{RGB}{180,30,30}
\definecolor{l3draw}{RGB}{40,80,130}
\definecolor{l3fill}{RGB}{210,220,235}
\definecolor{l4draw}{RGB}{30,65,110}
\definecolor{l4fill}{RGB}{200,212,228}
\definecolor{l5draw}{RGB}{100,35,0}
\definecolor{l5fill}{RGB}{230,210,200}
\definecolor{s1draw}{RGB}{75,85,95}
\definecolor{s1fill}{RGB}{215,218,222}
\definecolor{confborder}{RGB}{160,20,30}
\definecolor{conffill}{RGB}{245,210,210}

\newcommand{\cs}[1]{\texttt{\char`\\#1}}
\newcommand{\dataset}{\textsc{CTSpinoPelvic1K}}

% ── Visual-abstract figure: colours, styles, book-page macro, render slots ──
\definecolor{vacspine}{RGB}{120,42,0}
\definecolor{vacpelv}{RGB}{80,22,100}
\definecolor{vacfuse}{RGB}{25,100,180}
\definecolor{vactrain}{RGB}{18,78,40}
\definecolor{vacpseudo}{RGB}{175,95,0}
\definecolor{vacrib}{RGB}{150,30,30}
\definecolor{vacgray}{RGB}{90,95,105}
\tikzset{
  pback/.style={draw=#1!55!black, fill=#1!18, rounded corners=2pt,
                minimum width=1.85cm, minimum height=2.2cm, line width=0.4pt},
  pfront/.style={draw=#1!70!black, fill=#1!7, rounded corners=2pt,
                 minimum width=1.85cm, minimum height=2.2cm, line width=0.8pt,
                 inner sep=1.4pt},
  flow/.style={-{Stealth[length=2.6mm]}, line width=1.0pt, draw=vacgray},
  hdr/.style={font=\sffamily\bfseries\small, text=#1!55!black, align=center},
  % caption-below-a-render: category name then bold black count
  cap/.style={font=\sffamily\footnotesize, align=center, text=#1!55!black},
  lbl/.style={font=\sffamily\small, align=center, text=#1!50!black},
  procbox/.style={draw=#1!70!black, fill=#1!9, rounded corners=5pt,
                  line width=1.0pt, align=center, font=\sffamily, inner sep=8pt},
}
\newcommand{\num}[1]{\textcolor{black}{\bfseries #1}}
% Three offset "pages"; #5 = page content (a placeholder now, an
% \includegraphics later). Front node is named #1 so arrows can attach.
\newcommand{\pagestack}[5]{%
  \node[pback=#4] at ($(#2,#3)+(0.22,0.22)$) {};
  \node[pback=#4] at ($(#2,#3)+(0.11,0.11)$) {};
  \node[pfront=#4] (#1) at (#2,#3) {#5};
}
% Placeholder shown until a real 2-D CT-with-mask overlay is supplied.
\newcommand{\vacph}[1]{\parbox[c][1.85cm][c]{1.55cm}{\centering
  \scriptsize\itshape\color{vacgray}#1\\[3pt]\tiny[2-D render]}}
% ── RENDER SLOTS — drop in a real overlay by redefining the slot, e.g. put
%   \renewcommand{\vacSpine}{\includegraphics[width=1.9cm,height=2.4cm]{images/va/spine.png}}
% just before \begin{figure*}. Paths are relative to the main .tex, so any
% figure in the Overleaf project works. Each render is a CT slice with its
% segmentation mask overlaid.
\newcommand{\vacSpine}     {\vacph{spine GT}}
\newcommand{\vacPelvis}    {\vacph{pelvic GT}}
\newcommand{\vacFused}     {\vacph{spine+pelvis}}
\newcommand{\vacSpineOnly} {\vacph{spine}}
\newcommand{\vacPelvicOnly}{\vacph{pelvis}}
\newcommand{\vacDense}     {\vacph{spine+pelvis}}
\newcommand{\vacRibs}      {\vacph{+ ribs}}

\jmlrvolume{[VOLUME \# TBD]}
\jmlryear{2026}
\jmlrworkshop{Machine Learning for Healthcare (MLHC) 2026}

\title[\dataset{}]{\dataset{}: A Unified CT Benchmark for Lumbosacral Transitional Vertebra Segmentation}

\author{%
  \Name{Anonymous Author(s)}
  \Email{anonymous@institution.edu}\\
  \addr Anonymous Institution
}

\linenumbers
\begin{document}

\maketitle

\begin{abstract}
Machine learning models for spine segmentation face two compounding
blind spots: a \emph{modality gap} (predominance of MRI-derived
benchmarks despite CT being the clinical standard for surgical
planning) and an \emph{anomaly gap} (near-complete absence of
Lumbosacral Transitional Vertebra (LSTV) annotation in existing CT
datasets despite a 5--35\% population prevalence).
We introduce \dataset{}, a CT-native benchmark of 1{,}153
colonoscopy-protocol abdominal CT volumes across 802 patients
with a unified 10-class spinopelvic segmentation labelmap and
LSTV phenotype labels from two independent radiology teams.
The benchmark is assembled from CTSpine1K and CTPelvic1K---two
independently curated public datasets sharing a patient cohort but
using incompatible NIfTI coordinate conventions---via a
patient-anchored pipeline: an exact \texttt{PatientID} join
resolves identity, after which intrapatient bone-HU maximization
selects the source CT series for each annotation, with no spatial
registration and no metadata heuristics.
We further describe a \emph{partial-annotation training protocol}:
only 342 of the 1{,}153 cases carry fused spine-plus-pelvis ground
truth, but the remaining 811 spine-only and pelvic-only cases can
be incorporated via nnU-Net v2's \texttt{ignore\_label} mechanism,
a 3.37$\times$ increase in effective training data without biasing
the model toward background on unobserved structures.
We train an nnU-Net v2 segmenter under this protocol and use it to
densify the partially annotated cases out-of-fold: for a patient
whose spine and pelvis were traced on different acquisitions we keep
the spine-annotated scan as that patient's volume and model-complete
only its pelvis, and the few model-completed spines (on pelvic-only
patients) are manually reviewed.
The fidelity-critical structures---lumbar counting and the
L5/L6-versus-sacrum call---are thus always radiologist ground truth;
only the pelvis is model-completed, since it is large high-contrast
bone far less prone to vertebral miscounting.
We further characterize systematic 0\% cross-source agreement
between the two radiology teams at the sacralization boundary---not
annotation noise but a genuine definitional boundary between
morphological and vertebral-counting criteria---and release both
labels with explicit disagreement flags.
A zero-shot TotalSegmentator benchmark concretely quantifies the
anomaly gap: junction Dice drops 19.6 points (0.770 $\rightarrow$
0.574) from Normal to Any-LSTV cases, with lumbarization cases
degrading to 0.421 junction Dice owing to TotalSegmentator's
structural inability to emit an L6 class.
The \dataset{} dataset is publicly released on HuggingFace, and the
preprocessing pipeline, conversion scripts, nnU-Net training code,
and trained 5-fold checkpoints are all public.
\end{abstract}

\begin{keywords}
Spine Segmentation Benchmark, Partial-Annotation Training,
Lumbosacral Transitional Vertebra, Multi-Dataset Label Fusion,
Inter-rater Disagreement, TCIA, Surgical Planning, CT
\end{keywords}

\paragraph*{Data and Code Availability}
The \dataset{} dataset (images, unified masks, and the two-source
LSTV metadata) is derived from public TCIA collections under their
respective data use agreements and is publicly released at
\url{https://huggingface.co/datasets/anonymous-mlhc/CTSpinoPelvic1K}.
The preprocessing and TotalSegmentator-benchmark code, the nnU-Net
partial-annotation training pipeline, and the trained 5-fold
cross-validation checkpoints are all public and reachable from the
dataset README.

\paragraph*{IRB} All scans from de-identified public TCIA
repositories.

% ── FIGURE: Visual abstract ──────────────────────────────────────────────────
% To use real renders, redefine the slots just before this figure, e.g.:
%   \renewcommand{\vacSpine}{\includegraphics[width=1.9cm,height=2.4cm]{images/va/spine.png}}
\begin{figure*}[t]
\floatconts{fig:visual_abstract}
  {\caption{\textbf{How \dataset{} is built.}
    Two source mask sets---CTSpine1K spine GT and CTPelvic1K pelvic GT---are
    matched per patient by bone-HU placement. Masks traced on the \emph{same
    acquisition} are merged (brace) into a \textsc{fused} case; masks on
    \emph{different acquisitions} stay \textsc{separate} ($\rightarrow$) as
    spine-only and pelvic-only volumes. All configurations train an nnU-Net v2
    under the ignore-label protocol (all 1{,}153 cases). The fused and
    spine-only cases (782) are then pelvis-completed out-of-fold by that model;
    ribs are pseudolabelled with \textsc{TotalSegmentator} and added in v3. The
    page stacks are placeholders for 2-D CT-with-mask overlay renders; the count
    below each is the number of volumes in that category.}}
  {\resizebox{\textwidth}{!}{%
  \begin{tikzpicture}[x=1cm, y=1cm]
  \useasboundingbox (-1.7,-6.3) rectangle (20.9,6.6);   % padding -> margins + caption gap

  % ── phase headers (kept well above the stacks) ──
  \node[hdr=vacgray]   at (0,    5.9) {\textsc{Source masks}\\[-1pt]\footnotesize two TCIA datasets};
  \node[hdr=vacgray]   at (5,    5.9) {\textsc{Patient-anchored}\\[-1pt]\footnotesize match (bone-HU)};
  \node[hdr=vactrain]  at (9.6,  5.9) {\textsc{Partial-annotation}\\[-1pt]\footnotesize training};
  \node[hdr=vacpseudo] at (14.5, 5.9) {\textsc{Pseudolabel}\\[-1pt]\footnotesize complete the pelvis};
  \node[hdr=vacrib]    at (18.6, 5.9) {\textsc{+ Ribs}\\[-1pt]\footnotesize TotalSegmentator, v3};

  % ── (1) source mask stacks (titles BELOW each render) ──
  \pagestack{spine}{0}{2.7}{vacspine}{\vacSpine}
  \node[cap=vacspine] at (0,1.2) {CTSpine1K spine GT\\[2pt]\num{784} masks};
  \pagestack{pelv}{0}{-2.7}{vacpelv}{\vacPelvis}
  \node[cap=vacpelv] at (0,-4.2) {CTPelvic1K pelvic GT\\[2pt]\num{714} masks};

  % ── (2) match: } brace -> fused; -> arrows -> separate ──
  \pagestack{sonly}{5}{3.7}{vacspine}{\vacSpineOnly}
  \node[cap=vacspine] at (5,2.2) {spine-only\\[2pt]\num{440}};
  \pagestack{fused}{5}{0}{vacfuse}{\vacFused}
  \node[cap=vacfuse] at (5,-1.5) {fused (spine+pelvis)\\[2pt]\num{342}};
  \pagestack{ponly}{5}{-3.7}{vacpelv}{\vacPelvicOnly}
  \node[cap=vacpelv] at (5,-5.2) {pelvic-only\\[2pt]\num{371}};

  % the } merge, drawn as an explicit path (spike points at Fused)
  \draw[line width=0.8pt, draw=vacfuse!55!black] (spine.east) -- (2.4,2.7);
  \draw[line width=0.8pt, draw=vacfuse!55!black] (pelv.east)  -- (2.4,-2.7);
  \draw[line width=1.3pt, draw=vacfuse!75!black, rounded corners=2pt]
    (2.4,2.7) -- (2.7,2.7) -- (2.7,0.35) -- (3.05,0) -- (2.7,-0.35) -- (2.7,-2.7) -- (2.4,-2.7);
  \draw[flow, draw=vacfuse!75!black] (3.05,0) -- (fused.west);
  \node[lbl=vacfuse, font=\sffamily\scriptsize, fill=white, inner sep=1pt] at (1.9,0) {same\\acquisition};
  % separate cases: straight arrows (one label per arrow)
  \draw[flow] (spine.east) to[out=0,in=180] (sonly.west);
  \draw[flow] (pelv.east)  to[out=0,in=180] (ponly.west);
  \node[lbl=vacgray, font=\sffamily\scriptsize, fill=white, inner sep=1pt] at (2.5,4.05) {different\\acquisitions};
  \node[lbl=vacgray, font=\sffamily\scriptsize, fill=white, inner sep=1pt] at (2.5,-4.05) {different\\acquisitions};

  % ── (3) training box ──
  \node[procbox=vactrain, minimum width=2.4cm, minimum height=5.0cm] (train) at (9.6,0)
    {\bfseries nnU-Net v2\\[2pt]ignore-label\\training\\[7pt]
     \footnotesize\textbf{all 1{,}153}\\\footnotesize cases};
  \draw[flow] (fused.east) -- (train.west);
  \draw[flow] (sonly.east) to[out=0,in=152] ($(train.west)+(0,1.9)$);
  \draw[flow] (ponly.east) to[out=0,in=208] ($(train.west)+(0,-1.9)$);

  % ── (4) pseudolabel completion ──
  \pagestack{dense}{14.5}{0}{vacpseudo}{\vacDense}
  \node[cap=vacpseudo] at (14.5,-1.5) {dense spine\,+\,pelvis\\[2pt]\num{782}};
  \draw[flow, draw=vacpseudo!75!black] (train.east) -- (dense.west)
    node[midway, above=2pt, font=\sffamily\scriptsize, text=vacpseudo!60!black, align=center]
    {fused 342\\+\\spine-only 440\\\textbf{= 782}};
  \node[lbl=vacgray, font=\sffamily\scriptsize, align=center] at (12.0,-3.5)
    {pelvic-only 371\\(already real pelvic GT)};

  % ── (5) ribs (TotalSegmentator) ──
  \pagestack{ribs}{18.6}{0}{vacrib}{\vacRibs}
  \node[cap=vacrib] at (18.6,-1.5) {spine\,+\,pelvis\,+\,ribs\\[2pt]\num{782}};
  \draw[flow, draw=vacrib!75!black] (dense.east) -- (ribs.west)
    node[midway, above=2pt, font=\sffamily\scriptsize, text=vacrib!60!black, align=center]
    {TotalSeg.\\ribs};
  \end{tikzpicture}}}
\end{figure*}
% ── END FIGURE ───────────────────────────────────────────────────────────────

% =============================================================================
\section{Introduction}
\label{sec:intro}

The lumbosacral junction is among the most clinically consequential
regions of the axial skeleton.
Surgical procedures at this level---transforaminal lumbar interbody
fusion (TLIF), pedicle screw instrumentation, laminectomy, and
sacropelvic fixation---require preoperative planning at millimeter
accuracy, and errors at the lumbosacral level carry the highest
medicolegal burden in spine surgery~\citep{Carrino2011EffectImaging}.
Accurate automated segmentation of the lumbar vertebrae, sacrum, and
pelvic girdle supports pedicle screw trajectory optimization,
sacropelvic parameter measurement, intervertebral implant sizing,
and intraoperative neuronavigation registration.

\paragraph{The LSTV problem.}
A Lumbosacral Transitional Vertebra (LSTV) is a vertebra at the
L5--S1 junction exhibiting morphological characteristics of the
adjacent segment~\citep{Castellvi1984LumbosacralDefects}.
\emph{Sacralization} occurs when the fifth lumbar vertebra undergoes
bony fusion with the sacrum, reducing the mobile lumbar spine to
four segments.
\emph{Lumbarization} occurs when the first sacral segment detaches,
mimicking an independent sixth lumbar vertebra.
Prevalence estimates range from 5\% to 35\% depending on imaging
modality and classification criterion~\citep{Sekharappa2014LumbosacralSignificance}.
Wrong-level surgery attributable to transitional anatomy
misidentification occurs in an estimated 4.5--15\% of cases where
preoperative level counting relied on plain film or MRI
alone~\citep{Konin2010LumbosacralRelevance}.

\paragraph{The modality gap.}
The most widely used spine segmentation benchmarks---SPIDER,
Spine Generic, and TotalSpineSeg---are built entirely on
MRI~\citep{GraafSPIDERBenchmark,Cohen-Adad2021Open-accessManufacturers,Warszawer2025TotalSpineSeg:MRI}.
Castellvi morphology~\citep{Castellvi1984LumbosacralDefects} is defined
by cortical bridging and transverse process ossification, features
CT resolves at sub-millimeter resolution and MRI does not.
A model trained on MRI and deployed on preoperative CT acquires an
implicit domain shift that degrades lumbosacral boundary delineation
at the junction where clinical decisions are made.

\paragraph{The anomaly gap.}
CT-trained benchmarks including
CTSpine1K~\citep{Deng2021CTSpine1K:Tomography},
VerSe~\citep{Sekuboyina2020VerSe:Benchmark},
TotalSegmentator-CT~\citep{Wasserthal2023TotalSegmentator:Images}, and
CTPelvic1K~\citep{Liu2021DeepModels} do not report or control for
LSTV prevalence, do not stratify splits by transitional anatomy,
and make no attempt to evaluate per-type performance at the
lumbosacral junction.
No prior ML benchmark has curated LSTV in CT with dense spinopelvic
segmentation labels at scale, nor analyzed disagreement structure
between independent LSTV annotation operationalizations.

\paragraph{The annotation uncertainty problem.}
Even when LSTV labels are attempted, independent radiologists
applying different operationalizations of the same morphological
concept can systematically disagree at the lumbosacral boundary.
Existing datasets resolve this by publishing a single adjudicated
label, discarding the disagreement signal.
The medical machine learning community increasingly recognizes
that inter-rater disagreement is signal, not
noise~\citep{Dawid1979MaximumAlgorithm}:
\dataset{} applies this principle specifically to the LSTV
confusion zone, where disagreement reflects a genuine definitional
ambiguity rather than annotator error.

\paragraph{The partial-annotation opportunity.}
A systematic side-effect of fusing independently curated annotation
sets is that most cases are \emph{partially} annotated: the spine
mask is present without a matching pelvic mask, or vice versa.
Prior work typically discards these, retaining only the intersection
where both annotations exist.
For \dataset{}, this would collapse a 1{,}153-case benchmark to just
342 usable training cases---a two-thirds reduction that especially
penalizes rare subgroups.
We describe how nnU-Net v2's native \texttt{ignore\_label} mechanism
enables principled training across all 1{,}153 cases by rewriting
unobserved background voxels to the ignore class, eliciting zero loss
contribution from unobserved structures without biasing the model.

\paragraph{Contributions.}
\begin{enumerate}
  \item \textbf{\dataset{} dataset.} 1{,}153 abdominal CT volumes
    across 802 matched patients with unified 10-class spinopelvic
    masks and 53 LSTV-labeled volumes (from 33 distinct patients)
    spanning four phenotypes, released as a public HuggingFace
    dataset.
  \item \textbf{Re-anchoring of CTSpine1K and CTPelvic1K masks that
    shipped without scan linkage.} The released masks carry no
    \texttt{SeriesInstanceUID}, so each must be re-attached to one
    of its patient's TCIA scans: an exact \texttt{PatientID} join
    resolves identity, then intrapatient bone-HU maximization
    assigns each of the 1{,}498 masks to its source series
    deterministically, with no metadata heuristics and no spatial
    registration. We also characterize and correct a previously
    undocumented Y-axis coordinate-convention mismatch that would
    silently corrupt all fused masks.
  \item \textbf{Partial-annotation training protocol.}
    We describe an nnU-Net v2 \texttt{ignore\_label} protocol that
    enables a 3.37$\times$ effective training-set expansion on this
    multi-source setting by using all 1{,}153 cases rather than
    only the 342 fully fused ones, and use it to model-complete the
    missing pelvis out-of-fold (vertebral labels are never
    pseudolabelled). Trained-baseline evaluation is deferred to
    future work (Section~\ref{sec:discussion}).
  \item \textbf{LSTV labels from two independent radiology teams,
    with disagreement flags.} We retain both the CTPelvic1K
    morphological and the CTSpine1K counting-based LSTV label per
    case, add an independent radiologist Castellvi read of every
    LSTV case, and document 0\% cross-source sacralization
    agreement---not annotation noise but a definitional boundary
    between two legitimate criteria---making the confusion zone
    explicit and available for uncertainty-aware training.
  \item \textbf{Zero-shot TotalSegmentator benchmark quantifying
    the anomaly gap.} A widely used open-source multi-organ CT
    segmentation model~\citep{Wasserthal2023TotalSegmentator:Images}
    loses 19.6 junction-DSC points from Normal to Any-LSTV cases
    (0.770~$\rightarrow$~0.574), collapsing to 0.421 on
    lumbarization cases owing to the absence of an L6 output class.
\end{enumerate}

\subsection*{Generalizable Insights about Machine Learning in the
Context of Healthcare}

\textbf{(1) Partial-annotation training via ignore-label materially
expands the usable training set when fusing federated annotation
sources.} On \dataset{}, restricting training to fused cases would
discard 70\% of the annotated data; ignore-label training uses all
1{,}153 cases with principled supervision masking. The pattern
applies directly to any multi-source dataset where different teams
annotate different structures on the same underlying images.

\textbf{(2) Inter-rater disagreement at anatomical boundaries is a
recoverable signal, not noise.}
Two independent radiology teams applying legitimate but distinct
operationalizations of the same morphological concept achieve 0\%
cross-source agreement on sacralization. Datasets should retain
both annotations with source flags; downstream training and
evaluation can then exploit rather than discard this disagreement.

\textbf{(3) Coordinate-convention validation is a prerequisite for
multi-source medical image fusion.}
The Y-axis reflection between CTSpine1K and CTPelvic1K is
undetectable from voxel spacing alone and would have silently
corrupted all fused masks. Any pipeline combining independently
processed NIfTI datasets must explicitly validate
affine-derived world coordinates on representative landmarks
before training.

% =============================================================================
\section{Related Work}
\label{sec:related}

\paragraph{CT spine and pelvis benchmarks.}
CTSpine1K~\citep{Deng2021CTSpine1K:Tomography} annotated 1{,}005 multi-source
CTs with per-vertebra segmentation, and
VerSe~\citep{Sekuboyina2020VerSe:Benchmark} contributed 460 scans in challenge
format.
Neither reports LSTV prevalence or stratifies evaluation by
transitional anatomy, and many VerSe-derived scans are cropped to a
narrow spinal field of view, discarding the pelvic anatomy required
for spinopelvic parameter computation.
TotalSegmentator-CT~\citep{Wasserthal2023TotalSegmentator:Images} labels 104
structures across 1{,}204 CTs, but provides no LSTV classification and
reports neighboring-class confusion as a known failure mode at the
lumbosacral junction.
CTPelvic1K~\citep{Liu2021DeepModels} encodes LSTV qualitatively in
COLONOG mask filenames, which we leverage as a primary label source.

\paragraph{MRI benchmarks and the modality gap.}
SPIDER~\citep{GraafSPIDERBenchmark} provides 447 lumbar MRIs.
TotalSpineSeg~\citep{Warszawer2025TotalSpineSeg:MRI} is MRI-native;
sacrum labels were bootstrapped via CT-to-MRI rigid registration,
introducing residual domain noise at the lumbosacral boundary.

\paragraph{Multi-dataset label fusion and partial-annotation
training.}
\citet{Li2025LumbarMechanism} demonstrated combining CTSpine1K and
CTPelvic1K improves joint segmentation using a 3D residual
encoder-decoder with linear self-attention at the bottleneck and
cross-scale feature fusion in the decoder.
That work treats the two datasets as independent sources, provides
no LSTV-stratified evaluation, and does not resolve their
incompatible coordinate conventions.
On the methodological side, \citet{Wasserthal2023TotalSegmentator:Images}
and recent nnU-Net work~\citep{Isensee2024Nnu-netSegmentation} have
established that ignore-label training enables principled use of
partial-annotation cases at scale.
We are the first to describe this mechanism applied to the
spine+pelvis multi-source fusion setting with explicit
LSTV-stratified evaluation.

\paragraph{LSTV in the ML literature.}
\citet{Kwak2025AutomatedModel} evaluated a CNN on lumbar radiographs for
binary LSTV detection.
BUU-LSPINE~\citep{Klinwichit2023BUU-LSPINE:Detection} contributes 3{,}600 lumbar X-rays with
LSTV annotation, demonstrating clinical demand for automated LSTV
recognition that \dataset{} extends into 3D surgical planning CT.
No prior ML benchmark has curated LSTV in CT with dense spinopelvic
segmentation labels at scale.

\begin{table*}[t]
\floatconts{tab:dataset_comparison}
  {\caption{Comparison of spine, pelvis, and adjacent CT benchmarks.
    \checkmark~= available, ---~= not available,
    $\dagger$~= via cross-modal registration,
    $\ddagger$~= qualitative filename encoding only.
    \dataset{} is the only benchmark combining fused labels,
    LSTV annotation, two independent LSTV reads, and a
    registration-free fusion pipeline applicable to arbitrary TCIA
    collection pairs.}}
  {\resizebox{\textwidth}{!}{\small\begin{tabular}{lccccccccc}
  \toprule
  \bfseries Dataset & \bfseries Mod. & \bfseries $N$ &
  \bfseries Vertebral & \bfseries Pelvic & \bfseries Fused &
  \bfseries LSTV & \bfseries 2-Read &
  \bfseries Pathological Focus & \bfseries CT-native \\
  \midrule
  VerSe~\citep{Sekuboyina2020VerSe:Benchmark}
    & CT & 460 & \checkmark & --- & --- & --- & ---
    & Normal anatomy & \checkmark \\
  CTSpine1K~\citep{Deng2021CTSpine1K:Tomography}
    & CT & 1{,}005 & \checkmark & --- & --- & --- & ---
    & Normal anatomy & \checkmark \\
  CTPelvic1K~\citep{Liu2021DeepModels}
    & CT & 1{,}184 & partial & \checkmark & --- & $\ddagger$ & ---
    & Normal + metal artifacts & \checkmark \\
  TotalSegmentator~\citep{Wasserthal2023TotalSegmentator:Images}
    & CT & 1{,}204 & \checkmark & \checkmark & \checkmark & --- & ---
    & General multi-organ & \checkmark \\
  SPIDER~\citep{GraafSPIDERBenchmark}
    & MRI & 447 & \checkmark & --- & --- & --- & ---
    & Normal anatomy & --- \\
  TotalSpineSeg~\citep{Warszawer2025TotalSpineSeg:MRI}
    & MRI & 1{,}404 & \checkmark & $\dagger$ & $\dagger$ & --- & ---
    & Normal anatomy & --- \\
  \midrule
  BUU-LSPINE~\citep{Klinwichit2023BUU-LSPINE:Detection}
    & X-Ray & 3{,}600 & \checkmark & --- & --- & \checkmark & ---
    & LSTV & --- \\
  \midrule
  \textbf{CTSpinoPelvic1K} (ours)
    & CT & \textbf{1{,}153} & \checkmark & \checkmark & \checkmark
    & \checkmark & \checkmark
    & \textbf{LSTV / congenital variants} & \checkmark \\
  \bottomrule
  \end{tabular}}}
\end{table*}

% =============================================================================
\section{Dataset Construction: Patient-Anchored Fusion of CTSpine1K
and CTPelvic1K}
\label{sec:generalizable}

CTSpine1K and CTPelvic1K share a common COLONOG patient cohort but
were independently curated: masks are released as NIfTI files
without the \texttt{SeriesInstanceUID} linkage to the primary DICOM
data, and the two teams used different \texttt{dcm2niix}
invocations, producing silently incompatible coordinate conventions
(Section~\ref{sec:placement}).
Fusing these annotation sets therefore requires solving two
problems: (i)~\emph{identity}---determining which pelvic mask
belongs to which patient in which dataset---and
(ii)~\emph{alignment}---determining which specific CT acquisition
each mask was drawn from, in world coordinates.
We solve both without spatial registration or heuristic metadata
fallbacks.

\paragraph{Identity is trivial; alignment is the problem.}
Identity is an exact equality on the DICOM \texttt{PatientID} tag,
which is preserved in the metadata of both the source DICOM series
and the released mask NIfTI volumes; equal \texttt{PatientID} values
identify the same patient with algebraic certainty. This resolves
825/825 patients with no heuristic, and we do not claim it as a
contribution. The real difficulty is \emph{alignment}: no mask
carries a \texttt{SeriesInstanceUID}, and CTPelvic1K ships its
pelvic masks without their CTs entirely, so nothing directly links a
mask to a specific TCIA scan. With 825 patients spanning 3{,}451
TCIA series, the only available metadata proxies---a filename
SeriesNumber, the NIfTI affine geometry, the slice count---leave
most assignments merely probable: of the 1{,}498 masks, only 5
match a scan with certainty from metadata alone and 107 are fully
ambiguous.

\paragraph{Alignment via intrapatient bone-HU maximization.}
Restricting candidate series to the patient's own TCIA acquisitions
reduces the search space from a global index
($\sim$$10^5$--$10^6$ CT series) to
$|\mathcal{S}_p| \approx 4$ candidates. Within this bounded set,
bone-HU maximization (Eq.~\ref{eq:bonehu}) serves as a
metadata-free oracle for mask-to-series assignment---useful because
the DICOM metadata fields normally relied upon
(\texttt{SeriesDescription}, \texttt{BodyPartExamined}) are
frequently missing or mislabeled in TCIA releases.
It places every one of the 1{,}498 masks (784 spine, 714 pelvic) on
its source series deterministically, with zero alignment failures;
the distribution is shown in Section~\ref{sec:quality}.
Re-anchoring these orphaned, as-released-unusable masks into one
unified, ML-ready benchmark is the core construction contribution.

\paragraph{Partial-annotation opportunity.}
Because CTSpine1K and CTPelvic1K annotate overlapping but
non-identical patient subsets, the fused benchmark is
\emph{partially annotated}: 342 of 1{,}153 cases carry both spine
and pelvic ground truth, 440 are spine-only, and 371 are pelvic-only.
Prior work typically discards partial cases, retaining only the
intersection. We instead describe an nnU-Net v2
\texttt{ignore\_label} protocol for training across all 1{,}153
cases (Section~\ref{sec:ignore_label}), yielding a $3.37\times$
effective training-set expansion relative to fused-only baselines.

\paragraph{Scope of generalization.}
The patient-anchored pipeline is specific to the CTSpine1K/CTPelvic1K
fusion demonstrated here. The underlying primitives---\texttt{PatientID}-metadata
matching for identity resolution, and tissue-HU-based
intrapatient series selection for alignment---are applicable
whenever two annotation sets share a TCIA patient cohort, but we
do not claim empirical validation beyond the single dataset pair
reported here. We discuss candidate extensions in
Section~\ref{sec:discussion}.

% =============================================================================
\section{Source Data}
\label{sec:source}

\dataset{} assembles three interlocking public TCIA resources
sharing a common patient cohort.

\paragraph{COLONOG.}
The CT Colonography collection~\citep{Smith2015DataCOLONOGRAPHY} provides
paired supine and prone abdominal CTs from 825 subjects.
Protocol: 100--120~kVp, in-plane resolution 0.625--0.977~mm, slice
thickness 0.8--1.0~mm, soft-tissue kernel.
The paired acquisitions---absent from all prior spine benchmarks---
enable position-invariant training.

\paragraph{CTSpine1K.}
CTSpine1K~\citep{Deng2021CTSpine1K:Tomography} annotated 784 COLONOG cases with
per-vertebra segmentation using the VerSe label scheme
(IDs 20--25 encode L1--L6; ID 26 encodes sacrum).

\paragraph{CTPelvic1K.}
CTPelvic1K~\citep{Liu2021DeepModels} provides four-class pelvic
masks (sacrum, left hip, right hip) for 714 of the same COLONOG
patients.
Despite originating from the same acquisitions, different
\texttt{dcm2niix} invocations produced a systematic Y-axis
reflection between label spaces (Section~\ref{sec:methods}).
CTPelvic1K mask filenames encode patient UID, SeriesNumber, and
LSTV phenotype as a qualifier, providing the primary LSTV label
source for \dataset{}.

\begin{table}[t]
\floatconts{tab:scan_stats}
  {\caption{COLONOG scan statistics for \dataset{}.}}
  {\small\begin{tabular}{lc}
  \toprule
  \bfseries Property & \bfseries Value \\
  \midrule
  Patients total                  & 825 \\
  \quad With both positions       & 694 \\
  Spine masks (CTSpine1K)         & 784 \\
  Pelvic masks (CTPelvic1K)       & 714 \\
  Matched patients (\dataset{})   & 802 \\
  Matched volumes (\dataset{})    & 1{,}153 \\
  \quad Fused (spine + pelvis)    & 342 \\
  \quad Spine-only                & 440 \\
  \quad Pelvic-only               & 371 \\
  \bottomrule
  \end{tabular}}
\end{table}

% =============================================================================
\section{Methods}
\label{sec:methods}

\subsection{Phase 1: Patient-Anchored Identity Resolution}
\label{sec:patientdb}

Patient identity is carried explicitly by the DICOM \texttt{PatientID}
tag, which is preserved in the metadata of both the source TCIA DICOM
series and the released CTSpine1K and CTPelvic1K mask NIfTI volumes.
We establish identity by exact equality on this field: two records
belong to the same patient if and only if their \texttt{PatientID}
values are equal. The match is \emph{algebraically certain}---a
metadata equality, not a heuristic---and resolves \textbf{825/825
patients (100\%)}; we do not claim it as a contribution, as it merely
bounds the search to a single patient. PatientDB (schema in
Appendix~\ref{app:patientdb}) is keyed on this \texttt{PatientID}.
The substantive problem is the Phase~2 series alignment below: the
\texttt{SeriesInstanceUID} did not survive the DICOM-to-NIfTI
conversion, so each mask must still be placed onto one of its
patient's several scans.

\subsection{Phase 2: Intrapatient Exhaustive Series Assignment}
\label{sec:phase2}

For mask $m$ belonging to patient $p$, we identify its source TCIA
series by exhaustive intrapatient search.
COLONOG patients have $|\mathcal{S}_p| \approx 4$ CT-quality series,
making this tractable.
For each candidate series, we compute bone-HU coverage:
\begin{equation}
f_\text{bone}(m, s) =
  \frac{\sum_{v \in \mathcal{L}(m,s)}
        \mathbf{1}[I_s(v) > \tau_\text{bone}]}{|\mathcal{L}(m,s)|}
\quad \tau_\text{bone} = 200~\text{HU}
\label{eq:bonehu}
\end{equation}
and select $s^*(m) = \argmax_{s \in \mathcal{S}_p} f_\text{bone}(m,s)$.
The same criterion applies to both spine and pelvic masks, differing
only in the geometric alignment function $\mathcal{L}(m,s)$.
This eliminates metadata-dependent heuristics: a correctly matched
spine mask exhibits $f_\text{bone} \gg 0.5$; placement over a
soft-tissue or scout series yields $f_\text{bone} \approx 0$
(Figure~\ref{fig:bonehu_dist}).
Algorithm~\ref{alg:phase2} in Appendix~\ref{app:algorithm} provides
pseudocode.

\subsection{Affine Placement: The Y-Axis Reflection}
\label{sec:placement}

The DICOM standard mandates an LPS$+$ patient basis while NIfTI
natively encodes RAS$+$.
When independent teams convert the same DICOM source using different
toolchains, the resulting NIfTI affines carry silently inverted axes.
CTSpine1K encodes center voxels at world $y \approx -181$~mm;
CTPelvic1K at $y \approx +181$~mm for the same patients.
Given NIfTI affine $A$ and voxel extent $N_y$:
\begin{equation}
y_\text{ctr} = A_{1,3} + \tfrac{N_y}{2} \cdot A_{1,1},
\quad
\text{flip} = \bigl(y_\text{ctr}^\text{mask} \cdot
  y_\text{ctr}^\text{ref} < 0\bigr)
\end{equation}
When $\text{flip}=\text{true}$, the mask array is reversed along
axis~1 and the affine updated exactly.
Detection succeeds on all 714 pelvic masks; zero failures.
CTPelvic1K masks are additionally z-cropped sub-volumes; the
starting z-offset is recovered by cross-correlating CT bone-density
and mask occupancy profiles.

\subsection{Unified Label Map}
\label{sec:labels}

CTSpine1K VerSe labels and CTPelvic1K labels are remapped into a
unified 10-class scheme: 0\,=\,background; 1--5\,=\,L1--L5
(VerSe 20--24); 6\,=\,L6 / lumbarized S1 (VerSe 25; LSTV only);
7\,=\,sacrum (VerSe 26 / CTPelvic1K label 1); 8\,=\,left hip
(CTPelvic1K 2); 9\,=\,right hip (CTPelvic1K 3).
Class~6 is the first evaluation class in any public benchmark
reserved for the lumbarized transitional body.
At voxel-level conflicts at the L5/sacrum boundary, CTPelvic1K
takes precedence for sacral extent.

\subsection{Partial-Annotation Training Protocol}
\label{sec:ignore_label}

Fusion status is derived post-hoc: a case is \textsc{fused} when
$s^*(\text{spine}) = s^*(\text{pelvic})$, \textsc{spine-only} when
only a spine mask exists, and \textsc{pelvic-only} when only a
pelvic mask exists.
Training on fused cases alone would discard 811 of 1{,}153 volumes
(70\%), with a disproportionate penalty on rare LSTV subgroups.
We instead define a training protocol incorporating all 1{,}153
cases via nnU-Net v2's \texttt{ignore\_label} mechanism.

Concretely, for each case we rewrite the saved label volume as
follows:

\begin{center}
\begin{tabular}{@{}ll@{}}
\toprule
\textbf{Config} & \textbf{Label rewrite} \\
\midrule
\textsc{fused}        & identity (all 10 classes valid) \\
\textsc{spine-only}   & voxels$=0 \rightarrow 255$;
                        classes 1--6 retained \\
\textsc{pelvic-only}  & voxels$=0 \rightarrow 255$;
                        classes 7--9 retained \\
\bottomrule
\end{tabular}
\end{center}

The \texttt{ignore} class 255 is declared in the nnU-Net
\texttt{dataset.json}; voxels at this value contribute zero to both
the cross-entropy and the Dice loss, so the model is supervised
only on structures present in each case's ground truth.
For a \textsc{spine-only} case, the model learns ``here are L1--L6''
without being penalized on the pelvic region; symmetrically for
\textsc{pelvic-only}.
Spine classes receive supervision from all 342 fused plus 440
spine-only cases (782 total); pelvic classes from 342 fused plus
371 pelvic-only cases (713 total).
This is a principled use of federated annotation sources.
Per-class DSC, HD95, and mIOU are reported on the fused cases within
each split, which carry complete ground truth.
The nnU-Net v2 baseline is trained under this protocol (5-fold
cross-validation; the code and the trained checkpoints are public);
its quantitative evaluation on the fused test split, and the
fused-only versus all-cases ablation that isolates the ignore-label
contribution, follow at camera-ready
(Section~\ref{sec:discussion}).

\paragraph{Densifying the partial cases.}
We additionally use the trained model to \emph{densify} the released
labels so that every shipped volume carries a complete spine+pelvis
labelmap, and the completion is deliberately one-sided.
For a separate-mode patient---spine and pelvis traced on different
acquisitions---we keep the \emph{spine}-annotated scan as that
patient's representative volume and model-complete only its pelvis,
predicted out-of-fold by the 5-fold model so that no case is
densified by a network that trained on it.
The choice is principled: the lumbar count and the
L5/L6-versus-sacrum decision---where wrong-level error
originates---stay radiologist ground truth, while the pelvis (large,
high-contrast bone with no counting ambiguity) is the only
model-generated region, so completion cannot introduce a miscount or
a vertebral class-mix.
The exception is the few \textsc{pelvic-only} patients, for whom no
spine acquisition exists: there the spine is model-completed, and
\emph{every} such case is manually reviewed by a radiologist before
release.
The released completed branches (raw and expert-corrected) therefore
carry no ignore-label voxels, and the only model-generated labels are
pelves on spine-annotated scans plus the handful of reviewed
pelvic-only spines.

\begin{table}[t]
\floatconts{tab:splits}
  {\caption{\dataset{} composition and five-fold stratified splits
    (by LSTV phenotype and fusion status). Training uses all
    configurations via ignore-label; per-class metrics are reported
    on the fused cases, which carry complete ground truth.}}
  {\small\begin{tabular}{lr}
  \toprule
  \bfseries Breakdown & \bfseries Volumes \\
  \midrule
  \multicolumn{2}{@{}l}{\itshape By annotated region} \\
  \quad Fused (10 classes)      & 342 \\
  \quad Spine-only (ignore)     & 440 \\
  \quad Pelvic-only (ignore)    & 371 \\
  \midrule
  \multicolumn{2}{@{}l}{\itshape By split} \\
  \quad Train                   & 805 \\
  \quad Validation              & 174 \\
  \quad Test                    & 174 \\
  \midrule
  \textbf{Total volumes}        & \textbf{1{,}153} \\
  LSTV-positive (any config)    & 53 \\
  \bottomrule
  \end{tabular}}
\end{table}

% =============================================================================
\section{LSTV Labels and the Lumbosacral Confusion Zone}
\label{sec:lstv_arbitration}

\subsection{Two Independent LSTV Reads}

\textbf{CTPelvic1K filename qualifiers} encode a
\emph{morphological criterion}: the annotator inspects for cortical
bridging, transverse process hypertrophy, and pseudoarticulation
with the sacral ala, assigning one of four classes (\textsc{Normal},
\textsc{Lumbarization}, \textsc{Semi-sacralization},
\textsc{Sacralization}) based on Castellvi
morphology~\citep{Castellvi1984LumbosacralDefects}.
This is the canonical radiological definition.

\textbf{CTSpine1K vertebral heuristics} apply a
\emph{vertebral counting criterion}: VerSe label~25 present
$\Rightarrow$ lumbarization; VerSe label~24 absent $\Rightarrow$
sacralization; labels 20--24 present without 25 $\Rightarrow$ normal.
The two reads are independent: the vertebral count comes from
CTSpine1K (25 of the 33 LSTV patients are vertebral-source), the
pelvic morphological flag from CTPelvic1K (8 pelvic).

\textbf{Independent radiologist Castellvi read.}
Beyond the two inherited source labels, our radiologists have
independently Castellvi-classified \emph{every} LSTV case
(33 patients; 53 volumes, since a separate-mode patient is two
scans), including the ambiguous cross-source-disagreement cases,
recording the transitional type, left/right laterality (e.g.\
type~IV: type~IIa right, type~IIIa left), the non-rib-bearing
vertebra count, and per-case notes. This clinical read---released as
a phenotype table on the dataset branches and reported here as a
subset audit---is what makes the LSTV labels a deliberate clinical
characterization rather than an inherited heuristic.

Figure~\ref{fig:lstv_panel} illustrates representative cases for
each LSTV phenotype.

% ── FIGURE: LSTV 3D panel ────────────────────────────────────────────────────
\begin{figure*}[t]
\floatconts{fig:lstv_panel}
  {\caption{\textbf{Representative \dataset{} lumbosacral
    segmentation masks for each LSTV phenotype.}
    All renders are from fused cases displayed in anterior clinical
    view at native voxel resolution.
    Label colors follow the 10-class scheme: L1 (royal blue),
    L2 (sky blue), L3 (cyan-teal), L4 (green), L5 (yellow),
    L6 (orange), spine sacrum (red), pelvic sacrum (magenta).
    \textbf{(A)}~\textit{Normal anatomy} (COLONOG token 2):
    five lumbar vertebrae above the sacrum in standard configuration.
    \textbf{(B)}~\textit{Sacralization --- morphological criterion}
    (COLONOG token 32): Castellvi-type sacral fusion encoded in the
    CTPelvic1K filename qualifier; the CTSpine1K vertebral counting
    criterion independently labels this case \textsc{Normal},
    exemplifying the 0\% cross-source sacralization agreement.
    \textbf{(C)}~\textit{Sacralization --- vertebral-count criterion}
    (COLONOG token 107): VerSe label~24 absent; only four lumbar
    vertebrae identifiable.
    No CTPelvic1K morphological qualifier encodes this case as
    abnormal, again illustrating the definitional non-overlap.
    \textbf{(D)}~\textit{Lumbarization} (COLONOG token 167):
    VerSe label~25 (L6, orange) present as a discrete additional
    lumbar body between L5 (yellow) and the sacrum.}}
  {%
    \includegraphics[width=\textwidth]{images/LSTV_panel_4x1.jpg}%
  }
\end{figure*}

\subsection{The Confusion Zone}

Table~\ref{tab:lstv_agreement} shows cross-source agreement.
Sacralization shows \textbf{0\% cross-source agreement}
($n=28$ CTPelvic1K sacralization cases, zero overlap with CTSpine1K
sacralization cases).
This is not annotation error: a patient exhibiting Castellvi Type~II
morphology (bilateral pseudoarticulation satisfying the morphological
criterion) can simultaneously retain five independently countable
lumbar vertebral bodies if the bridging is partial.
Both annotators are correct within their own criterion; the case
sits precisely at the boundary of both classification systems
(panel~B of Figure~\ref{fig:lstv_panel}).

\begin{table}[t]
\floatconts{tab:lstv_agreement}
  {\caption{Cross-source LSTV agreement. The 0\% sacralization
    agreement quantifies the lumbosacral confusion zone.}}
  {\small\begin{tabular}{lrrrr}
  \toprule
  \bfseries LSTV Class & \bfseries $N$ & \bfseries \%
    & \bfseries Cross-val & \bfseries Agreement \\
  \midrule
  Normal              & 1{,}100 & 95.4 & 100\% & \checkmark \\
  Lumbarization       & 22 & 1.9 & n/a$^\dagger$ & --- \\
  Semi-sacralization  &  3 & 0.3 & 67\%$^\ddagger$ & partial \\
  Sacralization       & 28 & 2.4 & 0\%$^\S$ & \textbf{0\%} \\
  \midrule
  \textbf{Total LSTV} & 53 & 4.6 & --- & --- \\
  \textbf{Total}      & 1{,}153 & 100 & --- & --- \\
  \midrule
  \multicolumn{5}{l}{\footnotesize
    $^\dagger$ Pelvic protocol omits lumbarization body.} \\
  \multicolumn{5}{l}{\footnotesize
    $^\ddagger$ $n=3$; 1 case UNKNOWN.
    $^\S$ Definition mismatch; see text.} \\
  \bottomrule
  \end{tabular}}
\end{table}

Figure~\ref{fig:confusion_zone} illustrates the three anatomical
configurations that produce systematic disagreement: Type~A
(partial bridging with retained L5 body; Castellvi Type~II), Type~B
(asymmetric ossification), and Type~C (definitional ambiguity
between ``sacralized L5'' and ``lumbarized S1'' without MRI
correlation).

% ── FIGURE: Confusion Zone TikZ ──────────────────────────────────────────────
% DO NOT MODIFY THIS FIGURE
\begin{figure*}[t]
\floatconts{fig:confusion_zone}
  {\caption{\textbf{The lumbosacral confusion zone.}
    \textbf{Left}: normal L5--S1 anatomy with discrete junction.
    \textbf{Center}: the Type~A configuration that accounts for the
    majority of the 0\% sacralization agreement:
    Castellvi Type~II pseudoarticulation satisfies the morphological
    criterion (sacralization) while a discrete L5 body satisfies the
    counting criterion (normal).
    The transition zone (orange, dashed outline) is the voxel region
    at L5/S1 boundary whose class assignment is ambiguous.
    \textbf{Right}: conceptual disagreement rate as a function of
    cortical bridging extent.
    The confusion zone (red bars) is concentrated at intermediate
    bridging; full fusion and absent bridging produce reliable
    agreement.
    Exact values will be reported from the placed manifest.}}
  {%
  \begin{tikzpicture}[x=1.0cm, y=1.0cm,
    every node/.style={font=\small},
    arr/.style={-Stealth, thick, color=arrowgray},
  ]

  %% ── Panel A: Normal anatomy ──────────────────────────────────────────
  \node[font=\bfseries\small] at (1.7, 6.5) {Normal L5--S1};
  \node[rectangle, rounded corners=3pt,
        minimum width=2.4cm, minimum height=0.55cm,
        draw=l3draw, fill=l3fill, thick, align=center]
    at (1.7, 5.9) {L3};
  \node[rectangle, rounded corners=3pt,
        minimum width=2.4cm, minimum height=0.55cm,
        draw=l4draw, fill=l4fill, thick, align=center]
    at (1.7, 5.25) {L4};
  \node[rectangle, rounded corners=3pt,
        minimum width=2.4cm, minimum height=0.55cm,
        draw=l5draw, fill=l5fill, thick, align=center]
    at (1.7, 4.6) {L5};
  \draw[thick, arrowgray, dashed] (0.3, 4.27) -- (3.1, 4.27);
  \node[font=\tiny\itshape, text=arrowgray, anchor=west]
    at (3.15, 4.27) {L5--S1};
  \node[rectangle, rounded corners=3pt,
        minimum width=2.4cm, minimum height=0.55cm,
        draw=s1draw, fill=s1fill, thick, align=center]
    at (1.7, 3.95) {S1};
  \node[font=\tiny, text=agreem, anchor=west] at (0.15, 2.9)
    {\checkmark~Morph.: \textsc{Normal}};
  \node[font=\tiny, text=agreem, anchor=west] at (0.15, 2.55)
    {\checkmark~Count.: \textsc{Normal}};
  \draw[thick, agreem!70!black, rounded corners=2pt]
    (0.05, 2.4) rectangle (3.35, 3.1);

  %% ── Panel B: Type A confusion ────────────────────────────────────────
  \node[font=\bfseries\small, text=disagr] at (7.1, 6.5)
    {Type A: Partial Bridging};
  \node[rectangle, rounded corners=3pt,
        minimum width=2.4cm, minimum height=0.55cm,
        draw=l3draw, fill=l3fill, thick, align=center]
    at (7.1, 5.9) {L3};
  \node[rectangle, rounded corners=3pt,
        minimum width=2.4cm, minimum height=0.55cm,
        draw=l4draw, fill=l4fill, thick, align=center]
    at (7.1, 5.25) {L4};
  \node[rectangle, rounded corners=3pt,
        minimum width=2.4cm, minimum height=0.55cm,
        draw=l5draw, fill=l5fill, thick, align=center]
    at (7.1, 4.6) {L5 (countable)};
  \node[rectangle, rounded corners=4pt,
        minimum width=2.6cm, minimum height=0.65cm,
        draw=confborder, fill=conffill, very thick,
        align=center, dashed]
    at (7.1, 4.0)
    {\footnotesize pseudoarticulation\\[-2pt]
     \tiny (Castellvi Type~II)};
  \node[rectangle, rounded corners=3pt,
        minimum width=2.4cm, minimum height=0.55cm,
        draw=s1draw, fill=s1fill, thick, align=center]
    at (7.1, 3.35) {S1};
  \node[font=\tiny, text=lstvcol, anchor=west] at (5.55, 2.75)
    {\textbf{X}~Morph.: \textsc{Sacralization}};
  \node[font=\tiny, text=agreem, anchor=west] at (5.55, 2.40)
    {\checkmark~Count.: \textsc{Normal}};
  \draw[thick, disagr!70!black, rounded corners=2pt]
    (5.45, 2.25) rectangle (8.75, 2.95);
  \node[font=\tiny\bfseries, text=disagr, anchor=west] at (5.55, 2.05)
    {CONFUSION ZONE};

  %% ── Panel C: Disagreement bar chart ──────────────────────────────────
  \node[font=\bfseries\small] at (13.2, 6.5)
    {Disagreement by Bridging Extent};
  \node[font=\tiny\itshape, text=metalgray] at (13.2, 6.1)
    {(conceptual; exact values from manifest TBD)};
  \draw[thick, metalgray] (10.4, 2.0) -- (16.0, 2.0);
  \draw[thick, metalgray] (10.4, 2.0) -- (10.4, 5.5);
  \fill[agreem!60!white, draw=agreem!80!black, thin]
    (10.6, 2.0) rectangle (11.1, 2.4);
  \fill[agreem!60!white, draw=agreem!80!black, thin]
    (11.2, 2.0) rectangle (11.7, 2.6);
  \fill[disagr!60!white, draw=disagr!80!black, thin]
    (11.8, 2.0) rectangle (12.3, 3.5);
  \fill[disagr!60!white, draw=disagr!80!black, thin]
    (12.4, 2.0) rectangle (12.9, 4.8);
  \fill[disagr!60!white, draw=disagr!80!black, thin]
    (13.0, 2.0) rectangle (13.5, 5.2);
  \fill[disagr!60!white, draw=disagr!80!black, thin]
    (13.6, 2.0) rectangle (14.1, 4.1);
  \fill[disagr!60!white, draw=disagr!80!black, thin]
    (14.2, 2.0) rectangle (14.7, 2.8);
  \fill[agreem!60!white, draw=agreem!80!black, thin]
    (14.8, 2.0) rectangle (15.3, 2.5);
  \fill[agreem!60!white, draw=agreem!80!black, thin]
    (15.4, 2.0) rectangle (15.9, 2.2);
  \node[font=\tiny, anchor=north] at (10.85, 1.92) {none};
  \node[font=\tiny, anchor=north] at (13.25, 1.92) {partial};
  \node[font=\tiny, anchor=north] at (15.65, 1.92) {full};
  \node[font=\tiny, text=metalgray, rotate=90, anchor=south]
    at (10.2, 3.75) {\# disagreements};
  \node[font=\tiny, anchor=south] at (10.85, 1.72)
    {\footnotesize $\leftarrow$ bridging extent $\rightarrow$};
  \draw[thick, disagr] (11.75, 5.4) -- (14.75, 5.4);
  \draw[thick, disagr] (11.75, 5.3) -- (11.75, 5.4);
  \draw[thick, disagr] (14.75, 5.3) -- (14.75, 5.4);
  \node[font=\tiny\bfseries, text=disagr, anchor=south]
    at (13.25, 5.4) {confusion zone};
  \draw[thick, agreem] (10.55, 2.75) -- (11.75, 2.75);
  \node[font=\tiny, text=agreem, anchor=east] at (10.5, 2.75)
    {agree};
  \draw[thick, agreem] (14.75, 2.65) -- (15.95, 2.65);
  \node[font=\tiny, text=agreem, anchor=west] at (16.0, 2.65)
    {agree};
  \end{tikzpicture}
  }
\end{figure*}
% ── END FIGURE ───────────────────────────────────────────────────────────────

% =============================================================================
\section{Dataset Quality Assessment}
\label{sec:quality}

Bone fraction $f_\text{bone}$ (Eq.~\ref{eq:bonehu}) provides an
automated series-conflict detector.
Figure~\ref{fig:bonehu_dist} shows the distribution across all
1{,}498 placed masks.
Spine: mean 68.3\%, SD 11.2\%, min 15.2\%.
Pelvic: mean 66.2\%, SD 8.9\%, min 36.8\%.
All 784 spine and 714 pelvic masks place on a valid series
(1{,}498 total, zero alignment failures): every mask has at least one
intrapatient series clearing the bone-HU threshold by a wide margin
(Figure~\ref{fig:bonehu_dist}).
Zero spine masks had fewer than four identifiable lumbar labels;
9 had exactly four (all concordant with CTPelvic1K sacralization
or semi-sacralization qualifiers).

% ── FIGURE: Bone-HU histogram ────────────────────────────────────────────────
\begin{figure}[t]
\floatconts{fig:bonehu_dist}
  {\caption{\textbf{Bone-HU coverage distribution across 1{,}498
    placed masks.}
    Spine ($N=784$, blue) and pelvic ($N=714$, orange) both show
    strong bimodal separation between correctly matched
    soft-tissue-kernel CT series ($f_\text{bone}>0.5$; dominant right
    mode) and mismatched or scout-image series
    ($f_\text{bone}\approx 0$; small left mode).
    The gap between modes validates bone-HU maximization as a
    series-selection criterion without any metadata dependence.}}
  {%
    \includegraphics[width=\columnwidth]{images/fig1_bone_pct_histogram.png}%
  }
\end{figure}

% =============================================================================
\section{Zero-Shot Benchmark: TotalSegmentator}
\label{sec:results}

We evaluate TotalSegmentator---a widely used open-source multi-organ
CT segmentation model~\citep{Wasserthal2023TotalSegmentator:Images}---
as a zero-shot baseline on \dataset{}.
This isolates the \emph{anomaly gap}: the performance cost, at the
lumbosacral junction, of a model whose training distribution and
label vocabulary were not designed with LSTV in mind.

\subsection{Inference Setup}

Off-the-shelf inference with the pretrained \texttt{total}
task~\citep{Wasserthal2023TotalSegmentator:Images}, executed in
\emph{fast mode} (\texttt{--fast}), which resamples the input
volume to 3~mm isotropic resolution before prediction.
Inference was run on a single NVIDIA L40S GPU (46~GB) across the
full \dataset{} matched cohort ($N=921$ volumes completing
inference: 272 fused, 340 spine-only, 309 pelvic-native), of which
46 are LSTV-positive (26 sacralization, 16 lumbarization, 4 other).
We remap TotalSegmentator's native output labels
(\texttt{vertebrae\_L1}--\texttt{vertebrae\_L5}, \texttt{sacrum},
\texttt{hip\_left}, \texttt{hip\_right}) to our 10-class scheme.
TotalSegmentator does not emit an L6 class, so class~6 performance
is zero by construction; for lumbarization cases whose ground
truth contains an L6 body, this contributes a hard failure mode
that directly quantifies the clinical gap.
Because TotalSegmentator was not trained on any \dataset{} case,
we evaluate it zero-shot on the full matched cohort rather than
restricting to the 52-case test split; this materially increases
the statistical power of the LSTV-subgroup analysis
(Section~\ref{sec:lstv_subgroup}).
The 3~mm fast-mode setting trades approximately 1--2 DSC points of
headline performance~\citep{Wasserthal2023TotalSegmentator:Images} for
$\sim$20$\times$ wall-clock speedup; we report fast-mode numbers as
a realistic deployment proxy, noting that the LSTV performance
gap reported below is a property of TS's label vocabulary
(no L6 class) and training distribution, not of the resampling
resolution.

\paragraph{Scoring protocol.}
Every metric is scored only against the labels a case actually
carries; no metric is ever computed against an unobserved region.
Per-class DSC (Table~\ref{tab:results}) is therefore reported on the
272 fused cases alone, which have complete ground truth. The
subgroup junction metrics (Table~\ref{tab:lstv_ts}, over all 921
inferred cases) are scored per configuration against only the
present labels: a spine-only case scores L1--L6, a pelvic-only case
scores sacrum/hips, a fused case scores all nine; the junction DSC
and error rate count voxels only where ground truth exists, and
cases lacking both L5 and sacrum in the ground truth are excluded.
The scoring code is in the released repository.

\subsection{Per-Class Performance on the Full Fused Cohort}

Table~\ref{tab:results} reports per-class DSC on the full fused
cohort ($n=272$).

\begin{table}[t]
\floatconts{tab:results}
  {\caption{TotalSegmentator zero-shot per-class DSC (\%) on the
    full \dataset{} fused cohort ($n=272$), fast mode (3~mm).
    Junction DSC is evaluated over a 40~mm axial window centered
    on the L5/S1 boundary. Class~6 (L6) is zero by construction:
    TotalSegmentator does not emit an L6 class.}}
  {\small\begin{tabular}{lcccc}
  \toprule
  \bfseries Model & \bfseries L1--L5 & \bfseries L6
    & \bfseries Sacrum & \bfseries Jxn DSC \\
  \midrule
  TotalSegmentator & 85.4 & \phantom{0}0.0$^\star$ & 77.4 & 76.6 \\
  \midrule
  \multicolumn{5}{l}{\footnotesize
    $^\star$ TotalSegmentator does not emit an L6 class
    (structurally zero).}\\
  \multicolumn{5}{l}{\footnotesize
    L1--L5 is the mean of L1=84.7, L2=84.6, L3=85.7, L4=86.0,
    L5=85.9.}\\
  \bottomrule
  \end{tabular}}
\end{table}

TotalSegmentator achieves 85.4\% mean L1--L5 Dice and 77.4\%
sacrum Dice on the full fused cohort, but drops to 76.6\% at the
L5/S1 junction window---and, critically, collapses further on
LSTV-positive cases (Section~\ref{sec:lstv_subgroup}).
This baseline is well motivated: TotalSegmentator is a widely used
open-source multi-organ CT segmentation model, and its failure mode
on \dataset{}'s LSTV subgroup constitutes the concrete anomaly gap
our benchmark is designed to close.

\subsection{LSTV Subgroup Analysis}
\label{sec:lstv_subgroup}

Because TotalSegmentator is zero-shot on \dataset{}, every matched
case is a valid test case; we therefore evaluate its LSTV-subgroup
performance on the full matched cohort rather than the 52-case
fused test split.
This yields 46 LSTV-positive cases (26 sacralization,
16 lumbarization, 4 other) alongside 875 normal-anatomy cases---a
sample size sufficient to resolve phenotype-stratified performance
differences at the junction.
Table~\ref{tab:lstv_ts} reports the stratified analysis.

\begin{table}[t]
\floatconts{tab:lstv_ts}
  {\caption{TotalSegmentator zero-shot performance stratified by
    LSTV phenotype (fast mode, 3~mm).
    Jxn~DSC = 40~mm axial-window Dice at the L5/S1 boundary;
    ErrRate = per-voxel class-mismatch fraction within that window;
    TS\_L6 = count of cases in which TS emitted any L6 voxel
    (structurally zero by label vocabulary).
    The 19.6-point junction-DSC gap between Normal and Any-LSTV
    concretely quantifies the clinical cost of the anomaly gap.}}
  {\small\begin{tabular}{lrccc}
  \toprule
  \bfseries Subgroup & \bfseries $n$
    & \bfseries Jxn DSC & \bfseries ErrRate & \bfseries TS\_L6 \\
  \midrule
  Normal              & 875 & 0.770 & 0.330 & 0 \\
  Any LSTV            &  46 & 0.574 & 0.317 & 0 \\
  \quad Sacralization &  26 & 0.651 & 0.213 & 0 \\
  \quad Lumbarization &  16 & 0.421 & 0.526 & 0 \\
  \bottomrule
  \end{tabular}}
\end{table}

Three findings stand out.
First, the junction DSC gap between Normal (0.770) and Any-LSTV
(0.574) is 19.6 points---a clinically large disparity at the
exact anatomic level where wrong-level surgery risk is
concentrated.
Second, the disparity is driven asymmetrically by phenotype:
sacralization cases retain 0.651 junction DSC, whereas
lumbarization cases collapse to 0.421.
The gap is mechanistic rather than mysterious: every
lumbarization case carries a ground-truth L6 body that TS cannot
emit (\texttt{TS\_L6}${=}0$ across all 16 cases), so the L5/S1
window is guaranteed to contain a voxel-count proportion of
unrecoverable error.
Third, per-vertebra Dice degrades monotonically from L1 (58.6\%)
to L5 (37.2\%) within the Any-LSTV subgroup, confirming that the
loss of performance concentrates at the caudal lumbar levels
rather than being distributed uniformly across the spine.

Per-case qualitative predictions for the 9 LSTV-positive cases in
the fused test split, together with agreement-status
stratification, are provided in Appendix~\ref{app:lstv_cases}.

% =============================================================================
\section{Discussion}
\label{sec:discussion}

\paragraph{Partial-annotation training unlocks the full data scale.}
For federated annotation sources---the normal case when fusing
independently curated datasets---discarding partial cases imposes a
silent two-thirds data penalty at the \dataset{} scale.
Ignore-label training lets every annotated voxel contribute
supervision only where its class is observed.
The pattern generalizes beyond spine/pelvis fusion: any multi-source
dataset where different teams annotate different structures on the
same underlying images can apply this mechanism directly.

\paragraph{The confusion zone is a scientific contribution.}
The 0\% cross-source sacralization agreement documents, at scale,
that two legitimate expert-level annotation strategies for the same
anatomical feature produce non-overlapping classifications at the
Castellvi Type~I--II boundary.
No prior spine ML benchmark has reported this because no prior
benchmark retained both annotation sources with their labels intact.
The \texttt{lstv\_agreement} flag makes the ambiguity explicit,
enabling uncertainty-aware training and calibration benchmarking
on known-ambiguous cases.

\paragraph{The anomaly gap is empirically measurable.}
The zero-shot TotalSegmentator benchmark directly quantifies the
clinical cost of the anomaly gap: a widely used open-source
multi-organ CT segmentation model loses 19.6 junction-DSC points
from Normal to Any-LSTV cases, and 34.9 points on lumbarization
specifically (0.770 $\rightarrow$ 0.421).
Of the lumbarization-case error, a mechanically identifiable
fraction traces to a single architectural decision---the absence of
an L6 output class---illustrating how label vocabulary choices made
during benchmark design propagate as hard failure modes years
later in deployment.

\paragraph{Implications for surgical planning AI.}
Transitional anatomy is a documented driver of wrong-level spine
surgery and altered
instrumentation~\citep{Konin2010LumbosacralRelevance,Carrino2011EffectImaging},
so confident mislabeling at L5/S1 carries direct clinical cost.
A segmentation model that outputs high uncertainty at L5/S1 for a
Castellvi Type~II case is therefore more useful to a surgeon than one
that outputs a hard incorrect label with high confidence.
\dataset{} provides the annotation structure necessary to train
and evaluate such models.

\paragraph{The coordinate convention problem is general.}
The Y-axis reflection between CTSpine1K and CTPelvic1K is an
instance of a broad failure class: whenever independent teams
convert the same DICOM source using different toolchains, silently
inverted NIfTI axes result.
Any multi-source fusion pipeline must validate affine-derived world
coordinates for representative landmarks.

\paragraph{Limitations.}
\dataset{} derives from colonoscopy-protocol CT (soft-tissue kernel),
which may not generalize to the bone-kernel protocols used in some
surgical planning workflows; we keep this an explicit scope
restriction rather than claiming surgical-CT transfer.
LSTV is a stratified subgroup, not the primary contribution: 53 LSTV
volumes from 33 distinct patients (a separate-mode patient is two
scans). We therefore frame \dataset{} as an LSTV-aware spinopelvic
benchmark, and the headline findings---0\% cross-source agreement and
the anomaly gap---do not depend on large $N$. The LSTV labels are not
merely inherited: every LSTV case received an independent radiologist
Castellvi read (Section~\ref{sec:lstv_arbitration}), and the 0\%
sacralization agreement reflects a genuine criterion difference, not
annotator error. Semi-sacralization has limited representation
($n=3$).
TotalSegmentator was evaluated in fast mode (3~mm isotropic
resampling); full-resolution (1.5~mm) inference would be expected
to recover 1--2 DSC points on standard anatomy, but the
label-vocabulary-driven lumbarization failure mode is
resolution-invariant.
The nnU-Net v2 baseline is trained (5-fold cross-validation; it
produced the out-of-fold pelvic completion), and its code and
trained checkpoints are public; what remains for camera-ready is the
quantitative trained-baseline evaluation on the fused test split and
the fused-only versus all-cases ablation that directly tests the
ignore-label claim. Consequently that head-to-head comparison is
described here but not yet empirically reported.
The COLONOG cohort (colonoscopy-screening age, 45--75, US academic
centers) may not represent the full distribution of spine morphology
in surgical planning populations.

\paragraph{Future directions.}
Immediate next steps include: evaluating the trained nnU-Net v2
baseline under the partial-annotation protocol described in
Section~\ref{sec:ignore_label} on the fused test split; reporting
leak-safe out-of-fold Dice of the model-completed pelves---the only
model-generated labels---against CTPelvic1K's pelvic-only ground
truth; running a fused-only vs.\ all-cases ablation to isolate the
contribution of ignore-label training; and benchmarking additional
architectures---including the cross-scale attention architecture of
\citet{Li2025LumbarMechanism} and other recent spine-pelvis fusion
models---on the released benchmark.
Within the spine domain, further extensions include Castellvi
sub-type classification (Type~I--IV), bone-kernel protocol
inclusion, and validation on LSTV-enriched external cohorts with
independent re-reading.
The patient-anchored pipeline introduced here is also applicable
in principle to other TCIA collection pairs that share a patient
cohort; empirical validation on such extensions is beyond the scope
of the present work.

% ACKNOWLEDGEMENTS (camera-ready only)
% \acks{Many thanks to all collaborators and funders.}

\bibliography{bib/references}

\newpage
\appendix

% =============================================================================
\section{PatientDB Schema}
\label{app:patientdb}

PatientDB organizes all data using the DICOM \texttt{PatientID} as the
primary key, with five entities:
\texttt{PATIENT} (root entity indexed by \texttt{PatientID});
\texttt{TCIA\_SERIES} (CT acquisitions per patient);
\texttt{SPINE\_MASK} and \texttt{PELVIC\_MASK} (annotation entities
from CTSpine1K and CTPelvic1K respectively); and
\texttt{PLACED\_CASE} (output rows in the final manifest,
linking a mask to its selected series and fusion type).
The separation of identity (Phase~1) from alignment (Phase~2)
reduces a globally intractable matching problem to a locally
bounded one per patient, and generalizes immediately to any TCIA
collection pair sharing a patient cohort.

% =============================================================================
\section{Pipeline Construction (Detailed)}
\label{app:pipeline}

Figure~\ref{fig:pipeline} illustrates the patient-anchored
construction pipeline in full.
All styles in the TikZ diagram use pure named colors as style
arguments; no compound color (of the form \texttt{A!n}) is passed
as a parameterized argument.

% ── FIGURE: Pipeline TikZ (moved to appendix) ────────────────────────────────
% DO NOT MODIFY THIS FIGURE
\begin{figure*}[t]
\floatconts{fig:pipeline}
  {\caption{\textbf{Patient-anchored \dataset{} construction pipeline.}
    \textbf{(a)}~Phase~1: canonical UID matching builds PatientDB.
    \textbf{(b)}~Phase~2: bone-HU exhaustive search assigns each
    mask to its optimal series; outcomes merge into the manifest.
    \textbf{(c)}~The manifest feeds the 10-class schema and
    two-source LSTV summary.
    $\star$~=~bone-HU winner.}}
  {\centering%
\makebox[\textwidth][c]{%
\begin{tikzpicture}[x=1.0cm, y=1.0cm,
  every node/.style={font=\small},
  dsbox/.style={rectangle, rounded corners=4pt, minimum width=2.5cm,
                minimum height=0.95cm, draw=#1!70!black, fill=#1!12,
                thick, align=center, inner sep=4pt},
  mainbox/.style={rectangle, rounded corners=5pt, minimum width=3.8cm,
                  minimum height=1.0cm, draw=#1!80!black, fill=#1!15,
                  very thick, align=center, inner sep=4pt},
  widebin/.style={rectangle, rounded corners=3pt, minimum width=5.6cm,
                  minimum height=0.8cm, draw=#1!70!black, fill=#1!10,
                  thick, align=center, inner sep=4pt},
  smallbox/.style={rectangle, rounded corners=3pt, minimum width=2.2cm,
                   minimum height=0.75cm, draw=#1!80!black, fill=#1!13,
                   thick, align=center, inner sep=3pt},
  seriesbox/.style={rectangle, rounded corners=2pt, minimum width=2.8cm,
                    minimum height=0.62cm, draw=metalgray!60!white,
                    fill=metalgray!6!white, thin, align=left, inner sep=4pt},
  winnerbox/.style={rectangle, rounded corners=2pt, minimum width=2.8cm,
                    minimum height=0.62cm, draw=winner!80!black,
                    fill=winner!10!white, thick, align=left, inner sep=4pt},
  panellabel/.style={font=\bfseries\normalsize, text=white,
                     fill=metalgray!80!black, rounded corners=3pt,
                     inner sep=3pt, minimum width=0.6cm},
  seclabel/.style={font=\bfseries\footnotesize, text=metalgray!60!black},
  arr/.style={-Stealth, thick, color=arrowgray},
  arrthin/.style={-Stealth, thin, color=metalgray},
  arrgr/.style={-Stealth, thick, color=winner!80!black},
]
\fill[panelbg, rounded corners=6pt] (-0.3,-0.5) rectangle (4.8, 8.5);
\fill[panelbg, rounded corners=6pt] (5.1,-0.5) rectangle (12.2, 8.5);
\fill[panelbg, rounded corners=6pt] (12.5,-0.5) rectangle (16.8, 8.5);
\node[panellabel] at (0.15, 8.2) {(a)};
\node[panellabel] at (5.45, 8.2) {(b)};
\node[panellabel] at (12.85,8.2) {(c)};
\node[seclabel, anchor=west] at (0.55, 8.2)
  {\textsc{Phase 1: PatientDB}};
\node[seclabel, anchor=west, text width=5.8cm] at (5.85, 8.2)
  {\textsc{Phase 2: Intrapatient Exhaustive Search}};
\node[seclabel, anchor=west] at (13.25, 8.2)
  {\textsc{Output Manifest}};
\node[dsbox=colonog] (colonog) at (0.8, 6.7)
  {\textbf{\color{colonog!80!black}COLONOG}\\[-2pt]
   \scriptsize 825 patients\\[-2pt]\scriptsize paired CTs};
\node[dsbox=ctspine] (ctspine) at (2.2, 5.0)
  {\textbf{\color{ctspine!80!black}CTSpine1K}\\[-2pt]
   \scriptsize 784 spine segs};
\node[dsbox=ctpelv] (ctpelv) at (3.4, 3.3)
  {\textbf{\color{ctpelv!80!black}CTPelvic1K}\\[-2pt]
   \scriptsize 714 pelvic masks};
\node[rectangle, rounded corners=3pt, draw=metalgray!50!white,
      fill=white, minimum width=4.0cm, inner sep=4pt,
      align=center] (canon) at (1.8, 1.8)
  {\footnotesize\textbf{Canonical UID Matching}\\[1pt]
   \footnotesize $\text{canon}(u)=\text{prefix}\,\|\,
   \text{int}(\text{suffix})$};
\draw[arr] (colonog.south) -- ++(0,-0.45)
  -| ($(canon.north west)!0.15!(canon.north east)$);
\draw[arr] (ctspine.south) -- ++(0,-0.28)
  -| ($(canon.north west)!0.50!(canon.north east)$);
\draw[arr] (ctpelv.south)  -- ++(0,-0.18)
  -| ($(canon.north west)!0.85!(canon.north east)$);
\node[mainbox=patdb, minimum width=4.0cm,
      minimum height=1.3cm] (patdb) at (1.8, 0.15)
  {\textbf{\color{patdb!80!black}PatientDB}\\[1pt]
   \scriptsize \texttt{patient\_id}\,\textit{[PK]}\quad
   \texttt{tcia\_series[]}\quad\texttt{masks[]}};
\draw[arr] (canon.south) -- (patdb.north);
\node[widebin=patdb] (ptreg) at (8.6, 7.0)
  {\small Patient Registry\quad $p \in \{1,\ldots,802\}$};
\node[smallbox=ctspine] (smask) at (6.7, 5.7)
  {\footnotesize\textbf{Spine Seg} $m_s$\\[-2pt]
   \scriptsize world-space affine};
\node[smallbox=ctpelv] (pmask) at (10.5, 5.7)
  {\footnotesize\textbf{Pelvic Mask} $m_p$\\[-2pt]
   \scriptsize Y-flip + z-xcorr};
\draw[arr] (ptreg.south) -- ++(0,-0.2) -| (smask.north);
\draw[arr] (ptreg.south) -- ++(0,-0.2) -| (pmask.north);
\node[winnerbox] (sc1) at (6.7, 4.6)
  {\scriptsize $s_1$: prone\quad $f_\text{bone}{=}72\%$\;$\star$};
\node[seriesbox] (sc2) at (6.7, 3.9)
  {\scriptsize $s_2$: supine\quad
   $f_\text{bone}{=}\phantom{0}5\%$};
\node[seriesbox] (sc3) at (6.7, 3.2)
  {\scriptsize $s_3$: scout\quad
   $f_\text{bone}{=}\phantom{0}0\%$};
\draw[arrthin] (smask.south)--(sc1.north);
\draw[arrthin] (sc1.south)--(sc2.north);
\draw[arrthin] (sc2.south)--(sc3.north);
\node[rectangle, rounded corners=2pt, draw=winner!70!black,
      fill=winner!8!white, inner sep=3pt, minimum width=2.8cm,
      align=center] (sarg) at (6.7, 2.3)
  {\scriptsize $s^*(m_s){=}\displaystyle\argmax_{s}f_\text{bone}$};
\draw[arrgr] (sc3.south)--(sarg.north);
\node[seriesbox] (pc1) at (10.5, 4.6)
  {\scriptsize $s_1$: prone\quad
   $f_\text{bone}{=}\phantom{0}8\%$};
\node[winnerbox] (pc2) at (10.5, 3.9)
  {\scriptsize $s_2$: supine\quad $f_\text{bone}{=}71\%$\;$\star$};
\node[seriesbox] (pc3) at (10.5, 3.2)
  {\scriptsize $s_3$: scout\quad
   $f_\text{bone}{=}\phantom{0}0\%$};
\draw[arrthin] (pmask.south)--(pc1.north);
\draw[arrthin] (pc1.south)--(pc2.north);
\draw[arrthin] (pc2.south)--(pc3.north);
\node[rectangle, rounded corners=2pt, draw=winner!70!black,
      fill=winner!8!white, inner sep=3pt, minimum width=2.8cm,
      align=center] (parg) at (10.5, 2.3)
  {\scriptsize $s^*(m_p){=}\displaystyle\argmax_{s}f_\text{bone}$};
\draw[arrgr] (pc3.south)--(parg.north);
\node[diamond, draw=arrowgray, fill=white, thick,
      minimum width=2.2cm, minimum height=1.6cm,
      inner sep=2pt, align=center] (fuse) at (8.6, 1.25)
  {\footnotesize$s^*_s{=}s^*_p$?};
\draw[arr] (sarg.south) -- ++(0,-0.15) -| (fuse.west);
\draw[arr] (parg.south) -- ++(0,-0.15) -| (fuse.east);
\node[smallbox=fuseblue] (fused) at (6.7, 0.1)
  {\footnotesize\textsc{fused}\\[-2pt]\scriptsize same CT};
\node[smallbox=sepamber] (sep) at (10.5, 0.1)
  {\footnotesize\textsc{separate}\\[-2pt]\scriptsize own CT each};
\draw[arr] (fuse.west)
  -- node[above,font=\footnotesize,fill=white,inner sep=1pt,pos=0.3]
     {\textit{yes}}
  (fused.north east);
\draw[arr] (fuse.east)
  -- node[above,font=\footnotesize,fill=white,inner sep=1pt,pos=0.3]
     {\textit{no}}
  (sep.north west);
\node[rectangle, rounded corners=5pt, draw=metalgray!60!black,
      fill=white, very thick, minimum width=3.8cm,
      minimum height=2.4cm, inner sep=6pt,
      align=left] (mbox) at (14.65, 1.5)
  {\footnotesize\texttt{placed\_manifest.json}\\[3pt]
   \scriptsize\textcolor{fuseblue}{\rule{6pt}{6pt}}\enspace
     \textsc{fused}: $N=342$\\[1pt]
   \scriptsize\textcolor{sepamber}{\rule{6pt}{6pt}}\enspace
     \textsc{separate}: $N=351$\\[1pt]
   \scriptsize\textcolor{ctspine!70!black}{\rule{6pt}{6pt}}\enspace
     \textsc{spine\_only}: $N=89$\\[1pt]
   \scriptsize\textcolor{ctpelv!70!black}{\rule{6pt}{6pt}}\enspace
     \textsc{pelvic\_only}: $N=20$\\[2pt]
   \scriptsize\textbf{1{,}498/1{,}498 placed}};
\node[rectangle, rounded corners=4pt, draw=metalgray!50!white,
      fill=white, thick, text width=2.9cm,
      align=center, inner sep=3pt] (schema) at (14.65, 4.6)
  {\footnotesize\textbf{10-Class Schema}\\[3pt]
   \resizebox{2.6cm}{!}{%
   \begin{tikzpicture}[baseline=0pt, inner sep=0pt,
                        x=0.38cm, y=0.34cm]
     \foreach \i/\col/\lbl in {
       0/colonog!70!white/{L1},
       1/colonog!50!white/{L2},
       2/colonog!35!white/{L3},
       3/ctspine!50!white/{L4},
       4/ctspine!35!white/{L5},
       5/yellow!55!orange/{L6},
       6/red!60!black/{Sc},
       7/orange!80!white/{LH},
       8/orange!50!white/{RH}}
     {
       \fill[\col, draw=metalgray!30!white, very thin]
         (\i, 0) rectangle (\i+0.9, 1);
       \node[font=\tiny] at (\i+0.45, 0.5) {\lbl};
     }
   \end{tikzpicture}}\\[2pt]
   \tiny L1-5\,(1-5)\enspace L6\,(6)\\[-2pt]
   \tiny Sc\,(7)\enspace LH\,(8)\enspace RH\,(9)};
\node[rectangle, rounded corners=3pt,
      draw=lstvcol!60!orange, fill=lstvcol!8!white,
      thick, inner sep=5pt, minimum width=3.2cm,
      align=center] (lstv) at (14.65, 7.0)
  {\scriptsize\textbf{53 LSTV cases}\\[-2pt]
   \tiny 22 lumb.\,$|$\,3 semi\,$|$\,28 sacral.\\[-2pt]
   \tiny \textbf{two-source flags included}};
\coordinate (junc) at (8.6, -0.72);
\draw (fused.south) -- ++(0,-0.45) -| (junc);
\draw (sep.south)   -- ++(0,-0.45) -| (junc);
\draw[arr] (junc) -- ++(0,-0.15) -| (mbox.south);
\draw[arrthin] (mbox.north) -- (schema.south);
\draw[arrthin] (schema.north) -- (lstv.south);
\end{tikzpicture}
}
  }
\end{figure*}
% ── END FIGURE ───────────────────────────────────────────────────────────────

% =============================================================================
\section{Algorithms and Implementation}
\label{app:algorithm}

\begin{algorithm}[h]
\caption{Phase~2: Intrapatient Exhaustive Series Assignment}
\label{alg:phase2}
\begin{algorithmic}[1]
\REQUIRE PatientDB $\mathcal{D}$, alignment function
  $\textsc{Align}(m, s)$ returning placed voxels $\mathcal{L}$
\REQUIRE Bone-HU threshold $\tau = 200$~HU
\ENSURE Placement manifest $\mathcal{M}$
\FOR{each patient $p \in \mathcal{D}$}
  \STATE $\mathcal{S}_p \leftarrow$ CT-quality series for patient $p$
  \FOR{each mask $m \in
    \{\textsc{SpineMasks}(p) \cup \textsc{PelvicMasks}(p)\}$}
    \STATE $f^* \leftarrow 0$, $s^* \leftarrow \varnothing$
    \FOR{each series $s \in \mathcal{S}_p$}
      \STATE $\mathcal{L} \leftarrow \textsc{Align}(m, s)$
      \STATE $f \leftarrow
        |\{v \in \mathcal{L} : I_s(v) > \tau\}| / |\mathcal{L}|$
      \IF{$f > f^*$}
        \STATE $f^* \leftarrow f$,\; $s^* \leftarrow s$
      \ENDIF
    \ENDFOR
    \STATE $\mathcal{M}[p][m] \leftarrow (s^*,\, f^*)$
    \STATE Persist aligned mask $\mathcal{L}^*$ to \texttt{placed/}
  \ENDFOR
\ENDFOR
\RETURN $\mathcal{M}$
\end{algorithmic}
\end{algorithm}

\paragraph{NIfTI orientation convention.}
All placed masks are saved in PIR (Posterior-Inferior-Right)
orientation, matching MONAI \texttt{Orientationd(axcodes=`PIR')}
at training time.

\paragraph{TotalSegmentator inference configuration.}
TotalSegmentator v2 (\texttt{total} task) was run in fast mode
(\texttt{--fast}), which triggers the 3~mm isotropic-resampling
model rather than the default 1.5~mm variant.
Inference was executed per case on a single NVIDIA L40S GPU
(46~GB); typical per-case wall-clock time was 20--60 seconds,
dominated by the two-stage resample/predict loop.
Of the 1{,}153 matched volumes, 921 completed inference
successfully and are reported in Tables~\ref{tab:results} and
\ref{tab:lstv_ts}; failures were attributable to input-header
edge cases and are excluded rather than imputed.
Output labels were remapped one-to-one onto the 10-class
\dataset{} scheme, with TS \texttt{sacrum} $\rightarrow$ class~7,
\texttt{vertebrae\_L1}--\texttt{L5} $\rightarrow$ classes 1--5,
\texttt{hip\_left}/\texttt{hip\_right} $\rightarrow$ classes 8/9.
Class~6 (L6) has no TS source class by construction.

\paragraph{Partial-annotation label rewrite.}
The conversion script \texttt{tools/convert\_hf\_to\_nnunet.py}
reads each case's label NIfTI; for
\textsc{spine-only} and \textsc{pelvic-only} cases, voxels with
value 0 are rewritten to 255 (the declared \texttt{ignore\_label})
before being written to \texttt{labelsTr/}.
\textsc{fused} cases are copied verbatim. The resulting
\texttt{dataset.json} declares:
\texttt{"labels":\{"background":0,\dots,"right\_hip":9,"ignore":255\}}.

\paragraph{nnU-Net configuration.}
The trained baseline uses nnU-Net v2
with the ResEncM planner, GPU memory target 100~GB, overwritten
plans name \texttt{nnUNetResEncUNetPlans\_100G}, configuration
\texttt{3d\_fullres}, and a 100-epoch trainer variant.
Loss: \texttt{DC\_and\_CE\_loss} (nnU-Net default). Precision: bf16
mixed. One H200 GPU (140~GB HBM3e). Trained under 5-fold
cross-validation; the \texttt{splits\_final.json} is generated by
the conversion script to match our stratified split assignment
exactly, and the resulting checkpoints are released.

\paragraph{Two-source LSTV flag schema.}
Each manifest record includes:
\texttt{lstv\_pelvic} (CTPelvic1K morphological label);
\texttt{lstv\_vertebral} (CTSpine1K counting heuristic);
\texttt{lstv\_agreement} (Boolean: both sources agree on non-Normal
classification);
\texttt{lstv\_confusion\_zone} (Boolean: sources disagree on LSTV
status, marking a confusion-zone case).

% =============================================================================
\section{LSTV Case Details}
\label{app:lstv_cases}

The 9 LSTV-positive test cases are distributed as:
2 fused (1 lumbarization, 1 sacralization), 3 spine-only
(1 lumbarization, 2 sacralization), and 4 pelvic-native
(2 lumbarization, 2 sacralization).
Per-case qualitative predictions, agreement status, and junction
segmentation overlays are provided with the released dataset as
\texttt{lstv\_test\_qualitative.pdf}; the small cohort size
precludes reporting sensitivity/specificity with meaningful
confidence intervals here.

\end{document}